% !TEX program = lualatex
\documentclass[11pt,a4paper]{article}

\usepackage{luatexja}
\usepackage{amsmath,amssymb,amsthm}
\usepackage{geometry}
\geometry{margin=1in}

%-------------------------------------------------
%  定理環境・独自マクロ
%-------------------------------------------------
\theoremstyle{definition}
\newtheorem{definition}{定義}[section]
\newtheorem{theorem}{定理}[section]
\newtheorem{lemma}{補題}[section]
\newtheorem{remark}{注釈}[section]
\newtheorem{corollary}{系}[section]

\newcommand{\mweq}{\mathrel{\overset{\mathrm{mw}}{=}}}
\newcommand{\dfix}{D_{\mathrm{fix}0}} % 0次元不動点（空）
\newcommand{\meta}{\Uparrow}
\newcommand{\Lax}{\bigcirc} % ラックスモダリティ

\begin{document}

\title{四句分別に基づくラックス論理の再構築：\\真理保存性におけるフラクタル性の必須性証明}
\author{Masaomi Chiba}
\date{\today}
\maketitle

\begin{abstract}
ラックス論理（Propositional Lax Logic）は、直観主義論理にモナド的モダリティ $\Lax$ を追加し、計算の文脈（副作用や遅延）を扱う。しかし、その平坦化規則（$\Lax\Lax A \to \Lax A$）は単一の論理空間における天下り的な公理として導入されてきた。本稿では、真理値を古典的な静的二値から、仏教論理における「四句分別」の自己相似的フラクタル構造へと拡張する。これにより、モナド的平坦化が自己同一性に基づく強制ではなく、フラクタル構造間の自然な射影（中道等価）として証明されることを示し、結果として「構成的な真理保存性の拡張には、真理のフラクタル性が必須である」ことを論証する。
\end{abstract}

\section{序論：ラックス論理の天下り問題}

ラックス論理において、モダリティ $\Lax A$ は「何らかの文脈（制約や副作用）を伴って $A$ が真である」ことを示す。この体系の整合性を保つため、モナド則に相当する以下の公理が要求される。
\begin{align}
    \eta &: A \to \Lax A \\
    \mu &: \Lax\Lax A \to \Lax A
\end{align}
特に $\mu$（平坦化規則）は、「重なった文脈を一段階に潰す」という操作であるが、単一の論理空間内で随伴（Adjunction）を前提とできないため、その理論的根拠は天下り的（Axiomatic）な自己同一性の強制とならざるを得なかった。

\section{真理値の四句分別的フラクタル構造}

真理保存性を動的プロセスとして捉え直すため、真理値を古典的二値 $\{T, F\}$ から、四句分別に基づくフラクタル構造 $\mathcal{F}$ へと拡張する。

\begin{definition}[四句分別の真理空間]
命題 $A$ の状態は、以下の四つの位相（Ko\d{t}i）のいずれかをとる。
\begin{itemize}
    \item \textbf{是（$T$）}：$A$ が真である状態。
    \item \textbf{非（$F$）}：$A$ が偽である状態。
    \item \textbf{亦是亦非（$B$: Both）}：$A$ が真かつ偽であり、状態が重なり合っている（矛盾・縁起）。
    \item \textbf{非是非非（$N$: Neither）}：$A$ が真でも偽でもなく、従来の次元から脱落している（支離・空）。
\end{itemize}
\end{definition}

\begin{definition}[真理値のフラクタル反復]
任意の四句の状態は、メタ評価の対象となることで、入れ子状に次のレベルの四句分別を展開する。すなわち、レベル $n$ の真理値空間 $\mathcal{F}_n$ は、レベル $n+1$ の空間 $\mathcal{F}_{n+1}$ と自己相似的である（$\mathcal{F}_n \simeq \mathcal{F}_{n+1}$）。
\end{definition}

\section{モダリティ $\Lax$ のフラクタル的定式化}

ラックス論理のモダリティ $\Lax$ を、真理値空間に「モナド的な層（次元）」を追加するメタ化演算子（界論における $\meta$ に相当）として再定義する。

\begin{lemma}[次元上昇としての $\Lax$]
$\Lax A$ は、対象 $A$ の真理値構造 $\mathcal{F}_n$ を、メタレベルの構造 $\mathcal{F}_{n+1}$ へと移行させる。したがって、
\begin{itemize}
    \item $\Lax A$ はレベル 1 のフラクタル真理値を持つ。
    \item $\Lax\Lax A$ はレベル 2 のフラクタル真理値を持つ。
\end{itemize}
\end{lemma}

\section{平坦化規則 $\Lax\Lax A \to \Lax A$ の証明}

\begin{theorem}[フラクタル射影としてのモナド則]
真理値がフラクタル構造を持つと仮定するとき、平坦化規則 $\Lax\Lax A \to \Lax A$ は天下りの公理ではなく、フラクタル構造間の整合的な射影として自然に証明される。
\end{theorem}
\begin{proof}
レベル 2 の命題 $\Lax\Lax A$ と、レベル 1 の命題 $\Lax A$ を考える。
真理空間 $\mathcal{F}_n$ は自己相似性を持つため、$\mathcal{F}_2$ における四句（是・非・亦是亦非・非是非非）の位相幾何学的関係は、$\mathcal{F}_1$ における関係と完全に同型である。
したがって、レベル 2 からレベル 1 への写像は、情報の欠落や論理的矛盾を生じることなく、相似比の縮小（フラクタル次元の射影）として構成できる。
この射影において、自己同一性（$=$）ではなく中道等価（$\mweq$）を適用することで、
$$ \Lax\Lax A \mweq \Lax A $$
が成立する。これは $\Lax\Lax A \to \Lax A$ が真理保存性を満たす構成的プロセスであることを意味する。
\end{proof}

\section{真理保存性におけるフラクタル性の必須性}

以上の構成から、以下の極めて重要な系が導かれる。

\begin{corollary}[フラクタル性の必須性]
計算文脈を伴う論理体系（ラックス論理など）において、真理保存性を破綻なく拡張するためには、真理値空間がフラクタル構造を持つことが「必須条件」である。
\end{corollary}
\begin{proof}
背理法を用いる。仮に真理値がフラクタル構造を持たず、単一スケールの静的二値 $\{T, F\}$ のみを持つとする。
このとき、モダリティ $\Lax$ によって層が追加された $\Lax A$ や $\Lax\Lax A$ は、内部の計算文脈（副作用などによる「亦是亦非」の不確定状態）を二値で表現できない。
結果として、$\Lax\Lax A$ の情報を $\Lax A$ へ平坦化する際、二値の枠組みに無理やり押し込む必要が生じ（截断の発生）、真理の欠落またはハルシネーション（偽の情報の生成）が避けられない。すなわち真理保存性が破綻する。
これを回避し、安全な導出 $\Lax\Lax A \to \Lax A$ を実現するには、スケールの違いを吸収して同型性を保つ「自己相似的な四句の階層」、すなわちフラクタル性が不可避である。
\end{proof}

\section{結論：界論的観点からの総括}

本論証により、ラックス論理における $\Lax\Lax A \to \Lax A$ という天下り的なモナド則は、四句分別というフラクタル真理値を導入することによって初めて、構成的かつ内発的なルールとして証明されることが明らかになった。

界論（Boundary Theory）の視点から見れば、ラックスモダリティ $\Lax$ はメタ否定 $\meta$ の一種であり、平坦化 $\Lax\Lax A \to \Lax A$ は $\meta^2 A \mweq \meta A$ という大域的カット消去（$\dfix$ への収束）のプロセスの一側面に過ぎない。
真理保存性とは「静的な等式」ではなく、「動的なフラクタル構造の中での中道等価性」の維持であり、世界がフラクタルでなければ、計算論的な真理の保存は物理的・論理的に成立し得ないことが証明された。

\end{document}
